% Generated draft: requires mathematical and source review.
\documentclass[11pt]{article}
\usepackage[margin=25mm]{geometry}
\usepackage{fontspec}
\setmainfont{Latin Modern Roman}
\usepackage{amsmath,amssymb,mathtools,unicode-math}
\setmathfont{Latin Modern Math}
\usepackage{enumitem,needspace,xurl,hyperref}
\setlength{\parindent}{0pt}
\setlength{\parskip}{5pt}
\setlength{\emergencystretch}{3em}
\newcommand{\N}{\mathbb{N}}
\newcommand{\Z}{\mathbb{Z}}
\newcommand{\Q}{\mathbb{Q}}
\newcommand{\R}{\mathbb{R}}
\newcommand{\C}{\mathbb{C}}
\newcommand{\F}{\mathbb{F}}
\newcommand{\A}{\mathbb{A}}
\newcommand{\PP}{\mathbb{P}}
\newcommand{\E}{\mathbb{E}}
\newcommand{\eps}{\varepsilon}
\newcommand{\dd}{\,\mathrm{d}}
\begin{document}
\section*{UnsolvedMath ID: 30002769}
\textbf{Source problem code:} catalogue entry
\subsection{Definitions and mathematical statement}
Let \((X,J)\) be a compact complex manifold, where \(J\colon TX\to TX\) is its integrable complex structure on the underlying real tangent bundle. Let \(Z\) be a smooth complex projective variety with complex structure \(J_Z\). An algebraic foliation \(D\) has, on its regular locus \(U\subseteq Z\), an algebraic complex tangent subbundle \(D|_U\subseteq T_Z|_U\) whose local sections are closed under Lie brackets. The supplied record does not specify whether singularities of the foliation are allowed or which global convention for algebraic foliations is intended.

To make the notation \(J_f\) precise, use the following usual interpretation of a transverse embedding. A smooth embedding \(f\colon X\hookrightarrow U\) is transverse with complementary dimension if, for every \(x\in X\),
\[
 df_x(T_xX)\oplus (D_{f(x)})_{\R}=T_{f(x)}Z,
\]
where the tangent spaces in this equality are regarded as real vector spaces. Thus, if \(q\colon TZ|_U\to N_D:=TZ|_U/D\) is the quotient map, then
\[
 A_f:=q\circ df\colon TX\longrightarrow f^*N_D
\]
is an isomorphism of real vector bundles. The complex structure \(J_{N_D}\) induced by \(J_Z\) on the quotient defines
\[
 J_f:=A_f^{-1}\circ f^*J_{N_D}\circ A_f.
\]
In particular, \(\operatorname{rank}_{\C}N_D=\dim_{\C}X\). Here an isotopy through transverse embeddings means a smooth family \((f_t)_{t\in[0,1]}\) in which every \(f_t\) satisfies the stated embedding and transversality conditions.

Under these interpretations, the supplied question is: for every compact complex manifold \((X,J)\), do there exist a smooth complex projective variety \(Z\), an algebraic foliation \(D\) on \(Z\), and an isotopy class \(\alpha\) of transverse embeddings \(X\hookrightarrow Z\), such that
\[
 \exists f\in\alpha\qquad J_f=J?
\]
The equality is equality of endomorphisms of \(TX\). The class \(\alpha\) is existentially chosen; it is not prescribed in advance, and the statement requires equality for some representative, not every representative.

The source record does not define \(J_f\), specify complementary dimension, or specify the allowed isotopies and behavior at foliation singularities. The definitions above make explicit a conventional interpretation requiring the image to lie in the regular locus; these points require confirmation against the source before the draft is regarded as a definitive formulation.
\subsection{Short English statement}
The question asks whether every compact complex manifold can be smoothly embedded in a smooth complex projective variety so that the complex structure induced from the normal directions to an algebraic foliation is exactly its given complex structure. It also existentially specifies an isotopy class containing such an embedding. The precise conventions for transversality, isotopy, and foliation singularities remain to be confirmed from the source.
\subsection{Sources}
\begin{itemize}
\item Supplied local UnsolvedMath registry record, permanent ID \(30002769\), source problem code \texttt{catalogue entry}, title: \emph{Realizing Complex Manifolds Transversely to Algebraic Foliations}. The record belongs to the Oberwolfach Reports Open Problems collection and gives proposed year 2015. No individual proposer is supplied.
\item Supplied report citation: \emph{Geometric Methods of Complex Analysis} (2015), Oberwolfach Reports, DOI 10.4171/owr/2015/4, pp.\ 235--283. \url{https://doi.org/10.4171/owr/2015/4}. The supplied background describes these reports as documentation of mathematical workshops and their open-problem discussions; no problem-specific page is supplied.
\item Supplied literature reference: Demailly--Gaussier, JEMS. \url{https://doi.org/10.4171/JEMS/742}.
\item Supplied literature reference: Demailly--Gaussier, universal embedding spaces. \url{https://arxiv.org/abs/1905.06016}.
\item Supplied dated literature curation, checked at \texttt{2026-08-21}: the record labels the problem open and states that Demailly--Gaussier realize every compact almost complex manifold using an algebraic distribution, while the stronger Bogomolov question requiring an algebraic foliation on a projective variety remains posed as a basic question. This assessment is supplied dated curation derived from the OWR material, not an independently verified literature assessment. Its time-sensitive status and cited claims require independent verification.
\item The supplied statement status is \texttt{corrected\_verified}; its accompanying verification note reports only safe notation, typography, and TeX-delimiter cleanup, with hypotheses, quantifiers, and subquestions preserved. This registry label does not establish mathematical review of the definitions supplied in this draft.
\end{itemize}
\end{document}
